\keepXColumns{}
\begin{tabularx}{\textwidth}{
  >{\raggedright\arraybackslash}p{1cm} 
  >{\raggedright\arraybackslash}X
  p{4cm} 
  p{4cm} 
}
  \caption{Analisis Risiko dan Mitigasi}\label{tbl:risks} \\
  \toprule
  \textbf{Kode} & \textbf{Risiko} & \textbf{Dampak} & \textbf{Mitigasi} \\
  \midrule
  \endfirsthead

  \caption{Analisis Risiko dan Mitigasi (lanjutan)} \\
  \toprule
  \textbf{Kode} & \textbf{Risiko} & \textbf{Dampak} & \textbf{Mitigasi} \\
  \midrule
  \endhead

  \midrule
  \multicolumn{4}{l}{\textit{Bersambung ke halaman berikutnya}} \\
  %\bottomrule
  \endfoot

  \bottomrule
  \endlastfoot

  % === Table Data ===
  \hline
  R01 & Keterbatasan data audit & Analisis dan evaluasi sistem kurang representatif terhadap kondisi nyata gedung & Menggunakan data studi kasus yang tersedia dan melengkapi dengan data simulasi atau data historis yang relevan \\
  \hline
  R02 & Perubahan kebutuhan pengguna & Desain sistem perlu penyesuaian ulang sehingga menghambat implementasi & Membatasi ruang lingkup kebutuhan dan memprioritaskan kebutuhan inti sesuai rumusan masalah \\
  \hline
  R03 & Kendala teknis implementasi & Fitur tertentu tidak dapat diimplementasikan sesuai rencana & Menyederhanakan fitur nonkritis dan fokus pada fungsi utama audit energi \\
  \hline
  R04 & Keterbatasan waktu pengembangan & Implementasi dan pengujian tidak optimal & Menyusun jadwal kerja bertahap dan memfokuskan pengembangan pada modul utama sistem \\
  \hline
\end{tabularx}